\section{Introduction}
\label{sec:introduction}

Transformer encoders such as BERT remain attractive for text classification because they offer strong accuracy, stable training behaviour, and simple deployment interfaces. Their practical weakness is that a standard finetuned encoder is often much larger than the downstream task requires. The dense model spends compute on every attention head and every feed forward neuron at every layer, even when a large fraction of those structures can be removed without destroying the task solution. This report studies how to choose those structures under a strict inference budget and how to recover the accuracy of the resulting sparse architecture without retraining a new model from scratch.

The compression setting considered here is deliberately structured. The method removes complete attention heads and complete feed forward neurons, and in some cases removes an entire transformer layer by making it an identity map. This is different from unstructured weight pruning, where individual scalar weights are zeroed. The reason is operational rather than aesthetic. A dense smaller matrix is easy to export, run, profile, and compare; an irregular sparse matrix requires runtime support that may not exist on the deployment target and often fails to deliver a proportional speedup. The goal of the report is therefore not merely to reduce the number of non zero parameters, but to construct architectures whose forward pass is genuinely cheaper in the dimensions that matter for inference.

The central difficulty is that the architecture selection problem is combinatorial. A BERT-base encoder has twelve layers, twelve heads per layer, and 3072 feed forward neurons per layer. A candidate architecture is a binary decision over all of these structures. Exhaustive enumeration is impossible, and direct evaluation of a candidate by fully training or fully recovering it is too expensive to use inside a search loop. The method therefore separates the problem into two phases. The first phase searches over binary architectures using a cheap calibration metric. The second phase takes the selected architecture and performs a multi-stage recovery procedure based on distillation from the finetuned teacher.

The search phase is not a continuous relaxation of the masks. It is a direct search over hard binary masks. This choice avoids the gap between a relaxed gate and the binary architecture eventually deployed, but it also means that the search needs a robust sampling procedure. The report uses an evolutionary algorithm with a hard compute budget. The budget is expressed as a fixed operation count at a reference sequence length, and candidates are repaired into a narrow feasibility band around that budget. This constraint is essential: without it, a teacher-matching surrogate naturally favours keeping more structure, and the search degenerates into a size comparison rather than an allocation comparison. Inside the budget band, candidates have similar cost and differ mainly in where compute is placed across depth, attention, and feed forward capacity.

The surrogate metric used during search compares the masked student and the dense teacher on a calibration subset. It uses representation and prediction signals that are cheap enough to compute for many candidates. The report treats this metric as a practical instrument, not as an oracle. Several experiments show that the metric can mis-rank architectures, especially at high sparsity, because the unadapted masked student is not the same object as the recovered student. For that reason, the search includes diversity preservation, candidate caching, multi-fidelity evaluation, and a final funnel rather than trusting a single scalar score blindly.

The recovery phase is equally important. A selected mask can have very low raw accuracy before adaptation. The recovery procedure first hardens the masks and establishes the masked model state, then uses intermediate representation distillation to rebuild the internal trajectory of the student, and finally uses prediction distillation and task loss to improve the classifier-facing behaviour. The distinction between these stages matters empirically. On QNLI, for example, the representation recovery stage accounts for most of the accuracy restoration, while later prediction finetuning provides a smaller final correction. The report therefore describes the recovery procedure in detail rather than treating it as a generic finetuning step.

The contribution of the work is the full system: a hard-budget structured search over heads and neurons, a calibration metric designed for fast candidate evaluation, an evolutionary search procedure with diversity and caching machinery, and a multi-stage recovery method that makes the selected architecture usable. The experimental sections evaluate the resulting method on GLUE classification tasks, compare against simpler search baselines, ablate the search and recovery components, and diagnose where the metric succeeds or fails as a predictor of post-recovery accuracy.

The remainder of this chapter is organised as follows. Sections~\ref{chap:mask} and~\ref{chap:cost} define the candidate representation and cost model. Sections~\ref{chap:metric} through~\ref{chap:baselines} describe the surrogate metric and the search algorithm. Section~\ref{chap:recovery} describes the recovery procedure. The final sections report the experimental protocol and conclusions.
